\documentclass[border=10pt,crop=true]{standalone}
\usepackage{circuitikz}
\usetikzlibrary{calc, arrows.meta}
% Shared style for 电子学一 Chapter 2 op-amp circuit figures (CircuiTikZ).
\usetikzlibrary{calc}

\ctikzset{
  bipoles/length=0.95cm,
  bipoles/thickness=1,
  resistors/scale=1,
  capacitors/scale=1,
  label distance=2.5mm,
}

\tikzset{
  opfig/.style={font=\small},
  opfigThin/.style={font=\scriptsize},
}

% Geometry (cm) — keep identical across all op-amp schematics
\def\OpInLen{1.55}    % label → input terminal
\def\OpInGap{0.08}    % terminal → wire to op amp +
\def\OpAmpWire{1.00}  % terminal side → op amp + anchor
\def\OpOutLen{1.05}   % op amp out → output terminal
\def\OpJuncL{0.50}    % v_- → feedback junction (left)
\def\OpFbDown{0.55}   % v_- straight down before R_1
\def\OpRvert{1.20}    % R_1 vertical span to ground
\def\OpInSep{0.55}    % dual-input spacing (v_1 / v_2)
\def\OpFbStub{0.40}   % follower: stub right of output before dropping
\def\OpFbClear{0.22}  % follower: below op-amp south for horizontal feedback


% ---------------------------------------------------------------------------
% Spacing knobs — tweak here to adjust layout
% ---------------------------------------------------------------------------
\def\BusY{1.45}
\def\RfY{2.35}
\def\SepX{4.55}
\def\SepGap{0.11}
\def\SepH{2.55}
\def\RightX{5.55}         % right scope origin (clear of divider ~4.66)
\def\GndY{0}
% --- right model: isolated input loop vs output loop ---
\def\RxInBase{0.9}        % R_i / input node x (shift right past double line)
\def\RxInL{-0.55}         % i_i wire starts left of input node
\def\RxVSrcX{3.25}        % A i_i column (RxInBase + gap; no top tie to R_i)
\def\RxRoSpan{1.6}       % Vtop to R_o to output node horizontal span
\def\RxGndW{8.2}         % common ground rail length

\begin{document}
\begin{circuitikz}[american, scale=1.0, transform shape]

  % ========================================================================
  % LEFT — i-v converter
  % ========================================================================
  \begin{scope}
    \node[opfigThin, anchor=south] at (1.85, 2.95) {i--v converter};

    \coordinate (OAC) at (2.35, \GndY);
    \node[op amp, noinv input down](OA) at (OAC) {};

    \coordinate (VM) at (0.65, \BusY);
    \coordinate (Vo) at (3.55, \GndY);
    \coordinate (FBL) at (0.65, \RfY);
    \coordinate (FBR) at (2.55, \RfY);
    \coordinate (OutTap) at (2.55, \GndY);

    \draw (OA.+) -- ++(0, -0.95) node[ground, scale=0.85]{};
    \draw (OA.-) -| (VM);
    \node[circ, fill=black, inner sep=0.75pt] at (VM) {};

    \draw ($(VM)+(-1.05,0)$) to[I, i=$i_i$] (VM);
    \draw (VM) |- (FBL);
    \draw (FBL) to[R] (FBR);
    \node[opfigThin] at ($ (FBL)!0.5!(FBR)+(0,4mm) $) {$R$};
    \draw (FBR) -- (OutTap) -- (OA.out);
    \draw (OA.out) to[short, -o] (Vo) node[right=2.5mm, opfig]{$v_o$};
  \end{scope}

  % ========================================================================
  % MIDDLE — divider
  % ========================================================================
  \begin{scope}[xshift=\SepX cm]
    \coordinate (SepBot) at (0, \GndY-0.35);
    \coordinate (SepTop) at (0, \GndY-0.35+\SepH);
    \draw (SepBot) -- (SepTop);
    \draw ([xshift=\SepGap cm]SepBot) -- ([xshift=\SepGap cm]SepTop);
    \node[opfigThin, anchor=south west] at (0, 2.35) {model};
  \end{scope}

  % ========================================================================
  % RIGHT — port model
  %   Left:  i_i wire -> node -> R_i only (no top bus to the right)
  %   Right: A i_i (V) top -> R_o -> v_o / R_L  (isolated from R_i top)
  %   Bottom: single ground rail ties all shunt/source bottoms
  % ========================================================================
  \begin{scope}[xshift=\RightX cm]
    \coordinate (InL) at (\RxInBase+\RxInL, \BusY);
    \coordinate (In) at (\RxInBase, \BusY);
    \coordinate (GndIn) at (\RxInBase, \GndY);
    \coordinate (Vbot) at (\RxVSrcX, \GndY);
    \coordinate (Vtop) at (\RxVSrcX, \BusY);
    \coordinate (RoL) at (\RxVSrcX+\RxRoSpan, \BusY);
    \coordinate (RoR) at (\RxVSrcX+2*\RxRoSpan, \BusY);
    \coordinate (VoM) at (RoR);
    \coordinate (GndR) at (\RxVSrcX+2*\RxRoSpan, \GndY);
    \coordinate (GndEnd) at (\RxGndW, \GndY);

    % common ground (R_i, V source, R_L bottoms — no top connection between loops)
    \draw (GndIn) -- (GndEnd);
    \node[ground, scale=0.85] at ($(GndIn)!0.38!(GndEnd)$) {};

    % --- left input loop only ---
    \draw (InL) to[short, i=$i_i$] (In);
    \draw (In) to[R] (GndIn);
    \node[opfigThin, anchor=east] at ($ (In)!0.5!(GndIn)+(-5mm,0) $) {$R_i$};
    % (no \draw (In) -- (Vtop);  top buses intentionally disconnected)

    % --- right output loop only (V top connects solely to R_o left end) ---
    \draw (Vbot) to[V, l=$A i_i$] (Vtop);
    \draw (Vtop) -- (RoL);
    \draw (RoL) to[R] (RoR);
    \node[opfigThin] at ($ (RoL)!0.5!(RoR)+(0,4.5mm) $) {$R_o$};

    \draw (RoR) -- (VoM);
    \draw (VoM) to[short, -o] ++(0.55, 0) node[right, opfig]{$v_o$};
    \draw (VoM) to[R] (GndR);
    \node[opfigThin, anchor=west] at ($ (VoM)!0.5!(GndR)+(5mm,0) $) {$R_L$};
  \end{scope}

\end{circuitikz}
\end{document}
